%% 第 21 章：复谱定理
%% 本文件由 tools/md2tex.py 自动生成，请勿直接编辑。

\chapter{复谱定理}


\begin{jthm}{复谱定理}{r:22:1}
设 $V$ 是有限维复内积空间，$T:V\to V$。

下列条件等价：

\begin{jenum}
\item %
$T$ 是正规算子；
\item %
$V$ 存在一组由 $T$ 的特征向量构成的规范正交基；
\item %
$T$ 在某组规范正交基下具有对角矩阵。
\end{jenum}
\end{jthm}

\section{复谱定理的意义}

一般复算子只能保证上三角化，Jordan 块仍可能表示方向之间的耦合。

正规性保证这些耦合可以完全消失。空间可以由两两正交的特征方向组成，算子只在每个方向上独立地乘以一个复数特征值。

因此，正规算子在内积空间下是最理想的可对角化算子。

\begin{jproof}{复谱定理的证明思路}
Schur 定理保证：任意复内积空间上的算子，都能在某组规范正交基下表示为上三角矩阵。

设该上三角矩阵为 $A$。

若 $T$ 正规，则：

\[
\|Ae_1\|=\|A^*e_1\|.
\]

对于上三角矩阵，$Ae_1$ 只包含第一个对角元的贡献；而 $A^*e_1$ 还会包含第一行其余元素的贡献。

两个范数相等，只能说明第一行除对角元外的全部元素都为零。

于是 $A$ 被分成一个 $1\times1$ 块和一个更小的正规上三角块。对余下的块重复同样论证，最终所有非对角元素都为零。

因此，$A$ 必为对角矩阵。
\end{jproof}

\section{谱定理中的对角元}

若规范正交特征基为：

\[
(e_1,\ldots,e_n),
\]

且：

\[
Te_j=\lambda_je_j,
\]

则：

\[
[T]
=
\operatorname{diag}(\lambda_1,\ldots,\lambda_n).
\]

对角元就是对应基向量的特征值，按代数重数重复出现。

谱定理只能保证这种对角矩阵存在，但不给出每个特征值的具体数值，也就是说，谱定理保证对于这样的算子来说，存在某组正交基，且这组正交基的是整个空间的直和项，且每一项都是一维不变子空间，也就是一个特征向量。具体的特征值需要我们自己计算。

更一般的说法是：谱定理保证了这个算子在一个由它的特征向量构成的正交基下存在一个对角阵。

\section{规范正交特征基不唯一}

若某个特征值的特征空间维数大于 $1$，可以在该特征空间内部任取规范正交基。

不同特征值对应的特征空间彼此正交，因此将各特征空间中的规范正交基合并，仍得到整个空间的规范正交特征基。

即使全部特征值互异，也可以重排基向量；在复数域中，还可以分别给每个特征向量乘以模为 $1$ 的复数。

因此，谱定理通常给出许多不同的规范正交特征基。
